\documentclass[../../../include/open-logic-section]{subfiles}

\begin{document}

\olfileid{his}{set}{cantorplane}
\olsection{كانتور والخط والمستوى}

ترتبط بعض الملابسات المحيطة ببرهان شرودر--برنشتاين بالتاريخ الذي
ناقشناه في~\olref[his][set][pathology]{sec}. تذكر أن كانتور أثبت عام
1877 أن عدد نقاط المربع يساوي تمامًا عدد نقاط أحد أضلاعه. وسنعرض هنا
برهانه، أو محاولته الأولى للبرهان.

لتكن~$\unitline$ القطعة المستقيمة الواحدة، أي مجموعة النقاط~$[0,1]$.
وليكن~$\unitsquare$ مربع الوحدة، أي مجموعة النقاط~$\unitline \times
\unitline$. وبهذه المصطلحات، أثبت كانتور أن
$\cardeq{\unitline}{\unitsquare}$. وكتب إلى ديدكند ملاحظة تتضمن، في
جوهرها، الحجة الآتية.

\begin{thm}\ollabel{thm:cantorplane}
$\cardeq{\unitline}{\unitsquare}$
\end{thm}

\begin{proof}[البرهان: الجزء الأول.]
ثبّت~$a, b \in \unitline$. واكتبهما بالترميز الثنائي، فنحصل على
متتاليتين غير منتهيتين من الأصفار والآحاد،~$a_1$،~$a_2$،~\dots،
و$b_1$،~$b_2$،~\dots، بحيث:
\begin{align*}
a &= 0.a_1a_2a_3a_4\dots\\
b &= 0.b_1b_2b_3b_4\dots
\intertext{تأمل الآن الدالة~$f \colon \unitsquare \to \unitline$ المعطاة بـ}
f(a, b) & = 0.a_1b_1a_2b_2a_3b_3a_4b_4\dots
\end{align*}
تكون~$f$ !!a{injection}، لأنه إذا كان~$f(a, b) = f(c,d)$، كان
$a_n = c_n$ و$b_n = d_n$ لكل~$n \in \Nat$، ومن ثم~$a = c$ و$b = d$.
\end{proof}

لكن ذلك، كما بيّن ديدكند لكانتور، لا يجيب عن السؤال الأصلي. تأمل
$0.\dot{1}\dot{0} = 0.1010101010\ldots$. نحتاج إلى أن يكون
$f(a,b) = 0.\dot{1}\dot{0}$، حيث:
\begin{align*}
a&= 0.\dot{1}\dot{1} = 0.111111\ldots\\
b&= 0
\end{align*}
لكن~$a = 0.\dot{1}\dot{1} = 1$. ولذلك حين نقول «اكتب~$a$ و$b$
بالترميز الثنائي»، يجب أن نختار \emph{أي} الترميزين نستخدم؛ ولأن~$f$
ينبغي أن تكون \emph{دالة}، فلا يجوز لنا أن نستخدم إلا ترميزًا
\emph{واحدًا} من الترميزين الممكنين. لكن إذا استخدمنا مثلًا الترميز
المنتهي وكتبنا~$a$ على الصورة «$1.000\ldots$»، فلن يوجد زوج
$\tuple{a, b}$ بحيث~$f(a, b) = 0.\dot{1}\dot{0}$.

وخلاصة القول أن ديدكند أشار إلى أنه، بسبب إمكان وجود بعض التوسعات
الثنائية الدورية، تكون دالة كانتور~$f$ !!a{injection} لكنها
\emph{ليست} !!a{surjection}. وهكذا لم يبيّن كانتور إلا
$\cardle{\unitsquare}{\unitline}$، لا
$\cardeq{\unitsquare}{\unitline}$.

ورد كانتور على ديدكند على الفور تقريبًا، مقترحًا في جوهر كلامه أن
البرهان يمكن إتمامه كما يأتي:

\begin{proof}[البرهان: الإتمام.]
لقد بيّنا أن~$\cardle{\unitsquare}{\unitline}$. لكن من الواضح وجود
!!a{injection} من~$\unitline$ إلى~$\unitsquare$: ضع القطعة المستقيمة
منبسطة على أحد أضلاع المربع. ومن ثم
$\cardle{\unitline}{\unitsquare}$ و
$\cardle{\unitsquare}{\unitline}$. وبموجب شرودر--برنشتاين
(\olref[sfr][siz][sb]{thm:schroder-bernstein})،
$\cardeq{\unitline}{\unitsquare}$.
\end{proof}

لكن كانتور لم يكن يستطيع بالطبع إتمام السطر الأخير بهذه العبارات، إذ
لم تكن مبرهنة شرودر--برنشتاين قد بُرهنت بعد. والواقع أن كانتور، مع أنه
صاغها لاحقًا حدسًا عامًا، لم ير برهانًا مرضيًا لها حتى عام 1897.
(ولذلك قدم كانتور في وقت لاحق من عام 1877 برهانًا مختلفًا لـ
\olref{thm:cantorplane}، لم يمر عبر شرودر--برنشتاين.)

\end{document}
